\chapter{Process, Voltage and Temperature}
\label{ch:pvt}

No two chips are identical, and no chip runs in identical conditions all its life. Transistors come out of the factory a little faster or slower; the supply voltage sags under load; the temperature swings from a cold start to full load. A design must work across all of it. Engineers summarize these three sources of variation as \term{PVT}: process, voltage and temperature. This chapter explains each, finishes the temperature question from the assessment, and shows how designers decide which conditions to check.

\section{Process variation}

Manufacturing a chip involves printing features tens of nanometers wide and implanting dopant atoms by the thousands (Chapter~\ref{ch:fab}). Nothing can be done perfectly at that scale, so transistor properties vary:
\begin{itemize}
  \item \textbf{Global variation}: whole wafers or whole chips come out faster or slower than average (for example, a slightly thinner gate oxide everywhere on one wafer).
  \item \textbf{Local variation}: neighboring transistors on the same chip differ. With only a few dozen dopant atoms in a tiny transistor's channel, the random count alone changes its threshold voltage (\term{random dopant fluctuation}). The roughness of printed edges changes channel lengths.
\end{itemize}

Designers describe global variation with \term{process corners}, named by the speed of the NMOS and PMOS transistors:
\begin{center}
\begin{tabular}{ll}
\toprule
Corner & Meaning \\
\midrule
TT & typical NMOS, typical PMOS \\
FF & fast NMOS, fast PMOS \\
SS & slow NMOS, slow PMOS \\
FS & fast NMOS, slow PMOS \\
SF & slow NMOS, fast PMOS \\
\bottomrule
\end{tabular}
\end{center}
The skewed corners (FS and SF) matter for circuits that depend on a balance between pull-up and pull-down, such as an inverter's switching threshold or an SRAM cell's stability (Chapter~\ref{ch:memory}). Local variation is handled statistically, with \term{Monte Carlo} simulation (running many simulations with randomly varied transistors) and with timing margins called on-chip variation (OCV) derates.

A practical consequence is \term{binning}: chips from the same wafer are tested and sorted by how fast they run, and faster chips are sold as higher-end products.

\section{Voltage variation}

The supply voltage a transistor actually sees is lower than the regulator's output and changes over time (Chapters~\ref{ch:electricity} and~\ref{ch:capind}):
\begin{itemize}
  \item \textbf{Regulator tolerance}: the supply is specified as, say, $0.80\,\text{V} \pm 3\%$.
  \item \textbf{Static IR drop}: steady current through the resistance of the power grid lowers the voltage, more so far from the power pins.
  \item \textbf{Dynamic droop}: sudden increases in current cause $L\,di/dt$ dips until decoupling capacitors and the regulator respond.
\end{itemize}
Lower voltage means less overdrive, so transistors are slower. Power-delivery design aims to keep the total drop within a budget, often a few percent.

\section{Temperature}

Chips must work across a wide range of temperatures, from a cold start in an unheated room (or a car in winter) to full load under a heat sink. Typical junction-temperature ranges are about 0 to 100\,\textdegree C for consumer products, and $-40$ to 125\,\textdegree C or more for automotive parts. Temperature changes several things at once:

\begin{center}
\begin{tabular}{lll}
\toprule
Quantity & As temperature rises & Effect \\
\midrule
Carrier mobility $\mu$ & decreases & transistors slower \\
Threshold voltage $\Vt$ & decreases & transistors faster \\
Metal wire resistance & increases & wires slower \\
Subthreshold leakage & increases exponentially & more power, more heat \\
Electromigration & accelerates & shorter wire lifetime \\
\bottomrule
\end{tabular}
\end{center}

\subsection*{Why each happens}

\paragraph{Mobility falls.} Hot atoms vibrate more, so carriers collide with them more often (phonon scattering) and drift more slowly in the same electric field (Section~\ref{sec:tempmobility}).

\paragraph{Wire resistance rises.} Metal electrons are scattered by the same vibrations. Copper's resistance increases by about 0.39\% per \textdegree C.

\paragraph{Threshold voltage falls.} Heat gives carriers more energy, so a smaller gate voltage is enough to form the channel. A typical rate is on the order of 1\,mV per \textdegree C.

\begin{onoa}
The assessment asked how CMOS carrier mobility and wire resistance change as temperature increases. \textbf{Mobility decreases} and \textbf{wire resistance increases}; both make the circuit slower when hot. A complete answer also knows the exception below.
\end{onoa}

\begin{worked}{a hot wire}
A copper wire has a resistance of 100\,$\Omega$ at 25\,\textdegree C. What is its resistance at 85\,\textdegree C?
\[
  R = 100 \times [1 + 0.0039 \times (85 - 25)] = 100 \times 1.234 \approx 123\,\Omega.
\]
A 60\,\textdegree C rise adds about 23\% to the wire's resistance, and to its RC delay.
\end{worked}

\section{Temperature inversion}
\label{sec:tempinv}

For a transistor, heat has two opposite effects: lower mobility (slower) and lower threshold (faster). Which wins depends on the supply voltage. Drive current is roughly
\[
  I \propto \mu(T)\,\big(\VDD - \Vt(T)\big)^{\alpha}, \qquad \alpha \approx 1 \text{ to } 1.3.
\]

\begin{worked}{which effect wins? (illustrative numbers)}
Suppose heating a chip from 25\,\textdegree C to 100\,\textdegree C lowers mobility by 25\% and lowers $\Vt$ from 0.35\,V to 0.30\,V. Take $\alpha = 1.3$.

\textbf{At $\VDD = 1.0$\,V:} overdrive goes from 0.65 to 0.70\,V, up 7.7\%; raised to 1.3, about $+10\%$. Combined with the 25\% mobility loss: $0.75 \times 1.10 \approx 0.83$. Current \emph{falls}; the chip is slower when hot.

\textbf{At $\VDD = 0.5$\,V:} overdrive goes from 0.15 to 0.20\,V, up 33\%; raised to 1.3, about $+45\%$. Combined: $0.75 \times 1.45 \approx 1.09$. Current \emph{rises}; the chip is \emph{faster} when hot.
\end{worked}

At high supply voltages the overdrive is large, so a small change in $\Vt$ hardly matters and the mobility loss dominates: hotter is slower. At low supply voltages, common in modern power-saving chips, the overdrive is small, so the $\Vt$ drop is a large fraction of it and wins: hotter is faster, and the worst-case speed is at the \emph{coldest} temperature. This is \term{temperature inversion}.

\begin{trap}
True or false: ``A chip is always slowest at its highest temperature.'' False for modern low-voltage designs, because of temperature inversion. Wires, however, are always slower when hot.
\end{trap}

\section{Corners for sign-off}

Since no single condition is worst for everything, designers check a set of \term{PVT corners}:

\begin{center}
\begin{tabularx}{\linewidth}{>{\raggedright\arraybackslash}p{3.2cm}>{\raggedright\arraybackslash}X}
\toprule
Check & Worst-case corner \\
\midrule
Setup (paths too slow) & slow process (SS), low voltage, and the temperature that makes cells slowest (hot at high voltage; cold at low voltage, because of inversion) \\
Hold (paths too fast) & fast process (FF), high voltage, and the temperature that makes cells fastest \\
Leakage power & fast process (FF), high voltage, high temperature \\
Electromigration & high temperature, high current \\
\bottomrule
\end{tabularx}
\end{center}

\section{Aging}

Chips also slow down slowly over their lifetimes. Mechanisms such as NBTI (in PMOS) and hot-carrier injection gradually raise threshold voltages (Chapter~\ref{ch:fab}). Designers add an \term{aging guardband}: timing margin so that the chip still meets its frequency after years of use.

\section{Living with variation}

Modern chips measure their own conditions and adapt:
\begin{itemize}
  \item On-die \textbf{thermal sensors} trigger throttling before the chip overheats (Chapter~\ref{ch:heat}).
  \item \textbf{Ring oscillators} or replica critical paths measure how fast this particular chip is right now.
  \item \textbf{Adaptive voltage scaling} gives a fast chip a lower voltage (saving power) and a slow chip a higher one, so both meet the same frequency.
  \item \textbf{DVFS} picks voltage and frequency based on workload (Chapter~\ref{ch:power}).
\end{itemize}

\section{Summary}

\begin{itemize}
  \item Process: chips vary; corners TT, FF, SS, FS, SF; local variation handled statistically.
  \item Voltage: regulator tolerance, IR drop and $L\,di/dt$ droop lower the supply; lower voltage is slower.
  \item Temperature: mobility down, $\Vt$ down, wire resistance up, leakage up.
  \item At high $\VDD$ hot is slower; at low $\VDD$ hot can be faster (temperature inversion). Wires are always slower when hot.
  \item Setup is signed off at the slow corner, hold at the fast corner, leakage at fast, high-voltage, hot.
\end{itemize}

\section{Exercises}

\begin{enumerate}
  \item As temperature increases, what happens to carrier mobility, threshold voltage, wire resistance and leakage?
  \item A copper wire is 100\,$\Omega$ at 25\,\textdegree C. Estimate its resistance at 85\,\textdegree C.
  \item True or false: every chip is slowest at its maximum temperature.
  \item At which corner is setup checked? Hold? Worst-case leakage?
  \item A 0.80\,V supply suffers 5\% IR drop. What voltage do the transistors see, and are they faster or slower?
  \item What does the FS corner mean? Name a circuit property it might affect.
  \item Why is Monte Carlo simulation used?
\end{enumerate}

\subsection*{Answers}
\begin{enumerate}
  \item Mobility decreases; threshold voltage decreases; wire resistance increases; leakage increases (exponentially).
  \item About $100 \times (1 + 0.0039 \times 60) \approx 123\,\Omega$.
  \item False. At low supply voltage, temperature inversion makes cold the slow case.
  \item Setup: slow process, low voltage (and the slowest temperature). Hold: fast process, high voltage (and the fastest temperature). Leakage: fast process, high voltage, high temperature.
  \item 0.76\,V. Slower, because there is less overdrive.
  \item Fast NMOS, slow PMOS. It skews any balance between pull-up and pull-down, for example moving an inverter's switching threshold down, or affecting SRAM read and write margins.
  \item To estimate the effect of random local variation, by simulating many randomly varied copies of a circuit and looking at the spread of results.
\end{enumerate}
